\chapter{绪论}\label{chap:intro}

\section{研究背景及意义}

以大语言模型（Large Language Models，LLMs）为代表的生成式人工智能，已逐步应用于信息检索、内容创作、教育和智能服务等场景。模型从大规模语料中学习语言规律与世界知识，因而能够处理开放式任务；与此同时，训练数据、模型能力和用户输入也可能带来暴力、自残、歧视、违法指导等内容安全风险。越狱提示等攻击还可能诱导经过安全对齐的模型偏离预期行为\cite{zou2023gcg}。我国《生成式人工智能服务管理暂行办法》要求服务提供者承担相应的内容安全责任，并对违法内容采取处置措施\cite{cac2023measures}。在模型能力持续演进、部署方式日益多样的背景下，如何降低有害输出风险，同时保留模型对正常请求的服务能力，是大模型安全研究和应用中的重要问题。

目前，模型安全治理主要包括安全对齐、推理时检测与干预，以及对模型内部参数或知识的定向修改。安全对齐通常依赖监督微调、偏好优化等训练过程；推理时检测可以对输入或输出进行审核，但会增加系统组件和推理流程的复杂度\cite{ouyang2022instruct,rafailov2023dpo,inan2023llamaguard}。知识编辑则尝试定位与目标事实或行为相关的模型参数，并通过局部参数更新改变模型的特定输出倾向\cite{zhang2024comprehensive,meng2022rome,meng2023memit}。与全量再训练相比，这类方法有机会以较低成本完成发布后的定点干预。SafeEdit等研究已将知识编辑用于大模型安全治理，并提供了多类有害行为的评测基准\cite{wang2024safeedit}。需要指出的是，参数编辑所改变的是模型对目标请求的响应行为及其相关参数表征，不能仅凭输出变化就推断模型内部知识已被彻底删除；因此，编辑效果、泛化能力和副作用都需要通过系统评估加以检验。

面向安全场景，知识编辑至少面临两个相互关联的挑战。第一，编辑效果可能受到后续微调的影响。模型在部署与迭代过程中可能进行领域适配、指令更新或安全策略调整；已有研究在一般知识编辑任务中观察到后续微调会削弱部分编辑结果\cite{cheng2025finetune}。如果安全编辑随参数更新而衰减，原先受抑制的有害响应可能再次出现。第二，多类别安全问题具有不同的语义特征和风险等级，对模型行为的要求并不完全相同。统一编辑策略可能无法充分覆盖某些类别，也可能对与目标风险相近的正常请求产生过度干预。由此，安全知识编辑不仅需要关注编辑当下是否有效，也需要关注编辑能否经受后续微调、能否按风险类别和等级进行差异化控制。

围绕这些挑战开展研究具有理论与应用意义。在理论层面，本文研究面向微调鲁棒性的双子空间约束方法，分别考虑安全编辑目标与正常任务效用，并将后续微调扰动纳入优化过程，以拓展参数级安全编辑对持续更新环境的适应能力；同时研究多类别安全场景下的类别解耦与动态路由方法，探讨风险类别、等级与编辑模块之间的组织和调度关系。在应用层面，本文将相关方法集成到大语言模型生成内容安全增强系统中，支持编辑配置、风险路由和效果评估，为模型的发布后维护提供一种可控的参数级干预思路。

\begin{table}[htbp]
  \centering
  \caption{大模型内容安全治理技术路径比较}
  \label{tab:chap1-route}
  \begin{tabular}{p{0.14\textwidth}p{0.24\textwidth}p{0.27\textwidth}p{0.25\textwidth}}
    \hline
    技术路径 & 主要作用方式 & 优势 & 需要关注的问题 \\
    \hline
    安全对齐 & 通过训练数据和优化目标塑造模型行为 & 可在模型层面形成较广泛的安全行为 & 训练成本较高，后续微调可能改变对齐效果 \\
    推理时检测与干预 & 在输入、生成过程或输出阶段识别风险并采取措施 & 可在不改动基础模型参数的情况下追加策略 & 依赖额外组件与推理流程，检测和覆盖能力需评估 \\
    知识编辑 & 对与目标事实或行为相关的参数进行定点修改 & 适合发布后的局部更新，具有参数级干预潜力 & 需验证编辑持久性、泛化效果及对正常能力的影响 \\
    \hline
  \end{tabular}
\end{table}

\section{国内外研究现状}

\subsection{大模型生成内容安全}

大模型生成内容安全研究主要关注有害内容生成、偏见与歧视、隐私泄露以及越狱攻击等问题。Zou等人提出的通用对抗后缀方法展示了攻击提示绕过模型安全对齐的可能性\cite{zou2023gcg}。为使模型遵循人类指令与偏好，研究者提出了基于人类反馈的强化学习和直接偏好优化等方法\cite{ouyang2022instruct,rafailov2023dpo}。但安全对齐并不能保证模型在所有提示和部署条件下均保持安全；Qi等人的研究表明，少量有害样本的微调也可能损害已对齐模型的安全行为\cite{qi2023finetuning}。

机器遗忘将安全治理从输出行为约束进一步扩展到目标知识或能力的定向移除。WMDP基准围绕生物安全、网络安全等高风险知识评估模型能力，并提出表征误导向遗忘方法（Representation Misdirection for Unlearning，RMU）\cite{li2024wmdp}。负偏好优化（Negative Preference Optimization，NPO）从偏好优化角度构造遗忘目标，以缓解部分梯度上升方法带来的训练不稳定问题\cite{zhang2024npo}。约束知识遗忘（Constrained Knowledge Unlearning，CKU）则通过识别与有用知识相关的神经元并约束其梯度，在降低有害响应的同时关注通用能力保持\cite{shi2025cku}。这类方法说明，安全性与通用效用需要在同一优化过程中加以考虑；其主要特点是通过训练目标或梯度更新实现遗忘，和本文关注的参数级知识编辑在干预方式上有所不同。

知识编辑也逐渐被用于模型安全治理。Wang等人构建了SafeEdit基准，覆盖九类不安全行为，并提出基于神经元监测的去毒方法DINM，系统考察知识编辑在安全治理中的效果及其对通用能力的影响\cite{wang2024safeedit}。ToxEdit通过识别前向传播中的毒性激活模式并动态调整层间计算路径，关注毒性抑制与过度编辑之间的平衡\cite{lu2025toxedit}。DELMAN面向越狱防御，对与目标行为相关的参数进行动态编辑，并通过良性查询上的一致性约束降低对正常请求的影响\cite{wang2025delman}。上述研究表明，知识编辑已形成覆盖多种风险场景的探索，但不同工作针对的风险对象、编辑机制和评测范围并不相同，仍需进一步研究安全编辑的持久性及多类别场景下的协同控制。

\subsection{大模型知识编辑技术}

大语言模型在预训练过程中获得大量知识，但这些知识可能过时、错误或不符合特定应用中的安全要求。知识编辑旨在以局部干预的方式改变模型对指定事实或行为的响应，同时尽量保持无关知识和通用能力\cite{zhang2024comprehensive,wang2024survey}。按照干预位置与参数更新方式，相关方法大致可从外部知识补充、可训练编辑模块和定位后直接更新模型参数等路线加以理解。

（1）外部知识补充。检索增强生成（Retrieval-Augmented Generation，RAG）通过从外部知识库检索相关材料并将其提供给模型，降低知识更新对模型重训练的依赖\cite{lewis2020rag}。SERAC则维护编辑记忆，并通过分类器判断输入是否命中已有编辑，再选择编辑响应或原模型响应\cite{mitchell2022serac}。这类方法便于更新和撤销外部记录，但需要管理额外知识库或记忆模块，推理过程也依赖相应组件。对于安全治理而言，它们主要在模型外部补充或覆盖响应，模型参数本身并未因此发生定点修改。

（2）可训练编辑模块与轻量参数更新。LoRA通过低秩参数增量适配模型，减少可训练参数规模\cite{hu2022lora}；MEND利用元学习训练超网络，将编辑样本的梯度转换为参数修正量，以加快单条知识的更新\cite{mitchell2022mend}。此类方法可通过训练或附加参数改变模型行为，但需要预先训练编辑器或维护额外模块，其效果取决于训练数据、目标函数和模块配置。若编辑模块对应多个风险类别，如何控制模块之间的干扰以及推理时如何选择模块，也需要专门设计。

（3）定位后直接编辑模型参数。定位-编辑方法依据模型内部表征与目标知识的关联，确定候选参数位置，再计算局部权重更新。ROME利用因果追踪定位与单条事实相关的前馈层参数，并通过秩一更新写入目标事实\cite{meng2022rome}；MEMIT将编辑扩展到批量事实，并将更新分配到多个层\cite{meng2023memit}。AlphaEdit进一步将参数扰动投影到保留知识对应的零空间，以减少编辑过程对其他知识的影响\cite{fang2024alphaedit}。这些方法为参数级定点干预提供了重要基础，但现有方法的主要目标集中于事实更新与编辑过程中的知识保持，安全编辑在后续微调下的鲁棒性，以及面向多个风险类别的模块解耦与动态路由，仍需进一步研究。

在评测与工具方面，知识编辑研究通常考察目标编辑是否成功、相关提示上的泛化、无关事实的局部性以及文本流畅性等维度\cite{zhang2024comprehensive}。KnowEdit等基准和EasyEdit框架为多种编辑算法的复现、实现与比较提供了支持\cite{wang2024easyedit}。这些工作构成本文开展参数级安全编辑研究的技术基础。

\subsection{知识编辑持久性与鲁棒性}

参数级知识编辑通过局部权重更新改变模型输出，但一次编辑并不能自动保证其在后续参数更新后仍然有效。Cheng等人对多个模型编辑和微调配置进行了系统研究，发现后续微调可能造成已成功编辑的衰减，并指出微调参数范围及表征变化会影响编辑保持情况\cite{cheng2025finetune}。这一结果说明，编辑效果应放在模型的后续适配流程中考察，而不能只依据编辑完成时的静态指标进行判断。

与参数级编辑相关的另一条研究路线是微调鲁棒的机器遗忘。Wang等人将不变性约束引入大模型遗忘，使遗忘效果能在不同下游微调任务中得到考察\cite{wang2025ilu}；CKU则通过限制与有用知识相关的更新来兼顾安全性和通用能力\cite{shi2025cku}。这些方法为后续更新下的遗忘鲁棒性和效用保护提供了参考，但其核心操作仍建立在训练式遗忘或梯度优化过程之上。对于直接写入局部权重的知识编辑，如何在编辑时构造适当的更新约束，并面向未见的后续微调扰动优化编辑结果，仍是需要单独处理的问题。

在编辑过程中保护未修改知识方面，AlphaEdit通过零空间投影约束参数扰动，减少对保留知识的影响\cite{fang2024alphaedit}。该思路与本文关注的效用保护具有联系，但编辑阶段保护已有知识与编辑后抵抗持续微调是两个不同环节：前者约束本次参数更新，后者需要面对未来训练数据、优化方向和更新幅度的不确定性。因此，本文将安全编辑目标与正常任务效用分别纳入子空间约束，并进一步考虑后续微调造成的参数扰动，研究双子空间约束与鲁棒优化的结合方式。

现有研究由此留下三个具体问题：第一，普通事实编辑的保持率不能直接代表安全编辑失效后有害响应是否回潮，安全行为需要专门的持久性评估；第二，训练式遗忘的鲁棒化机制难以直接照搬到闭式或局部参数编辑，需要在编辑阶段设计适用的参数更新约束；第三，编辑阶段的知识保护不等同于微调后的安全保持，需在安全目标、正常任务效用和后续更新扰动之间建立统一的分析与评价框架。

\subsection{面向多类别风险控制的知识编辑研究现状}

安全风险包含多个类别，且相同类别中的请求也可能具有不同的危害程度和上下文。SafeEdit以九类不安全行为构建评测基准，为多类别安全编辑提供了数据和评价起点\cite{wang2024safeedit}。现有工作进一步探索了按风险特征调整编辑行为的可能性：ToxEdit根据毒性激活动态调整计算路径\cite{lu2025toxedit}，DELMAN围绕越狱行为实施动态参数编辑，并关注良性请求上的效用保持\cite{wang2025delman}。这些研究表明，动态干预和多风险评估已进入安全编辑研究视野。

不过，风险类别覆盖、类别模块组织和输入时路由属于不同层面的设计问题。一个基准包含多类测试样本，并不意味着编辑方法已经为各类别建立了相互独立的参数模块；能够检测毒性或越狱输入，也不等同于可以根据类别及风险等级选择不同编辑强度。对于联合编辑，不同类别可能共享部分参数表征，更新一个类别时可能影响其他类别或正常任务。因此，多类别场景仍需研究如何建立类别独立的低秩编辑模块、如何通过解耦约束减轻模块间干扰，以及如何依据输入风险特征动态调度模块。本文的第二项研究即围绕这些问题展开。

\subsection{现有方法存在的问题}

综合上述研究，面向大模型生成内容安全的知识编辑已在风险基准、局部参数干预和效用保持方面取得进展，但面向持续部署和多类别控制仍有以下问题需要研究。

（1）安全编辑在后续微调下的保持能力缺乏针对性研究。已有编辑持久性研究主要考察一般知识编辑在微调后的变化\cite{cheng2025finetune}；而对安全编辑而言，关注点还包括有害响应是否回潮、正常安全行为是否保持，以及不同微调目标造成的差异。需要建立适用于安全编辑的微调威胁模型和保持性评价方法。

（2）编辑阶段的安全目标和效用保持目标缺少统一的参数约束。零空间编辑等方法关注本次编辑对保留知识的影响\cite{fang2024alphaedit}，遗忘方法则通过训练目标或梯度约束兼顾安全与效用\cite{shi2025cku,wang2025ilu}。但对于参数级安全编辑，如何分别刻画安全编辑方向和正常任务效用方向，并将二者与后续微调扰动共同纳入优化，尚需进一步探索。

（3）多类别联合编辑中的类别干扰和风险差异尚未得到充分处理。已有安全编辑研究覆盖毒性、越狱等具体场景，SafeEdit也提供多类别评测\cite{wang2024safeedit,lu2025toxedit,wang2025delman}；然而，面向多类别联合编辑，仍需要考察编辑模块能否按类别隔离、类别间更新是否互相干扰，以及路由策略能否区分风险类别和等级。

（4）现有评价和系统组织尚未贯通编辑、持久性与动态路由。编辑成功率和正常任务表现不能单独说明编辑经过微调后的可靠性；总体安全分数也难以揭示某一类别的漏防或对某类正常请求的过度拒绝。因此，评价需要覆盖微调后编辑保持、类别级编辑效果、模块干扰、路由性能和正常请求效用，并将这些指标纳入统一的安全增强系统流程。

上述问题分别对应本文的两项方法研究和系统设计：第（1）（2）项构成微调鲁棒性研究的出发点；第（3）项构成多类别分级解耦与动态路由研究的出发点；第（4）项贯穿实验评价和系统实现。

\section{本文主要内容}

针对1.2.5节归纳的问题，本文以参数级知识编辑为技术主线，围绕安全编辑的微调鲁棒性、多类别风险控制和工程化应用开展研究。整体思路是先建立兼顾安全目标与正常任务效用的编辑约束，再面向多类别风险设计解耦编辑模块和动态路由机制，最后将相关方法组织为安全增强系统。本文主要研究内容如下。

（1）面向微调鲁棒性的双子空间约束安全知识编辑方法。针对安全编辑结果可能受到后续微调影响的问题，本文分别基于安全编辑目标和正常任务样本构建安全子空间与效用保护子空间，并研究如何在参数更新过程中同时利用这两个子空间约束编辑方向。在此基础上，对后续微调引起的参数扰动进行建模，将微调后的安全编辑效果纳入鲁棒优化目标。实验设计围绕编辑前后的安全行为、不同微调条件下的保持率、通用能力变化和编辑局部性展开，并通过消融与参数敏感性分析考察子空间约束和鲁棒优化各自的作用。该研究对应第二章。

（2）面向多类别安全性的分级解耦与动态路由编辑方法。针对多类别有害内容的编辑干扰和风险处置粒度不足问题，本文按照风险类别和等级组织类别独立的低秩安全编辑模块，并设计类别解耦约束以降低联合更新时的跨类别干扰；进一步利用输入内容的风险特征计算路由权重，动态选择或组合相应编辑模块。实验设计重点考察多类别编辑效果、类别隔离能力、风险等级识别与路由性能、正常请求效用及过度拒答，并通过消融和敏感性分析验证各模块的贡献。该研究对应第三章。

（3）大语言模型生成内容安全增强系统设计与开发。围绕上述两项方法，本文设计并开发安全增强系统，组织基础模型管理、鲁棒安全编辑、多类别安全编辑、风险感知动态路由、安全响应生成和安全性能评估等功能，形成从编辑配置、推理响应到结果评估的基本流程。系统章节将说明需求、架构、功能模块与测试方案，对方法的工程集成进行考察。该研究对应第四章。

\begin{table}[htbp]
  \centering
  \caption{本文研究内容与问题、方法及评价维度的对应关系}
  \label{tab:chap1-framework}
  \begin{tabular}{p{0.19\textwidth}p{0.23\textwidth}p{0.31\textwidth}p{0.16\textwidth}}
    \hline
    研究问题 & 研究方法 & 主要评价维度 & 对应章节 \\
    \hline
    编辑结果可能受到后续微调影响 & 安全子空间、效用保护子空间与扰动鲁棒优化 & 安全编辑效果、微调保持率、通用能力、局部性 & 第二章 \\
    多类别联合编辑可能相互干扰，风险处置粒度不足 & 类别独立编辑模块、解耦约束与风险感知动态路由 & 类别编辑效果、隔离性、路由性能、过度拒答 & 第三章 \\
    算法方法需要形成可操作的安全治理流程 & 集成编辑、路由、响应与评估功能的安全增强系统 & 模块集成与系统测试 & 第四章 \\
    \hline
  \end{tabular}
\end{table}

\section{本文章节安排}

本文共分为五章，各章安排如下。

第一章为绪论，介绍研究背景与意义，综述大模型生成内容安全、知识编辑技术、知识编辑持久性与鲁棒性，以及面向多类别风险控制的知识编辑研究；在此基础上归纳现有方法的问题，说明本文主要研究内容和章节安排。

第二章为面向微调鲁棒性的双子空间约束安全知识编辑方法。首先介绍参数级知识编辑、安全编辑问题及后续微调扰动；随后构建安全子空间和效用保护子空间，研究双子空间约束下的参数更新策略，并将后续微调扰动纳入鲁棒优化；最后通过对比、消融和敏感性实验，考察安全编辑效果、微调后保持能力及通用效用变化。

第三章为基于类别解耦与动态路由的分级安全知识编辑方法。首先说明多类别安全编辑、风险类别与等级建模等相关概念；随后设计类别独立的低秩编辑模块和类别解耦约束，并研究基于输入风险特征的动态路由及模块调度；最后通过多类别对比、类别隔离、路由性能、安全效用平衡和消融实验进行评价。

第四章为大语言模型生成内容安全增强系统设计与开发。该章分析系统需求，介绍系统概要、功能模块、开发实现与测试方案，重点说明鲁棒编辑、多类别编辑、动态路由和安全评估等功能的集成方式。

第五章为结论与建议。总结全文研究工作，结合方法和实验结果分析研究局限，并提出后续改进方向。
